\documentclass[10pt,letterpaper]{article}

\usepackage[letterpaper,top=0.43in,bottom=0.43in,left=0.48in,right=0.48in]{geometry}
\usepackage[T1]{fontenc}
\usepackage{newpxtext,newpxmath}
\usepackage{microtype}
\usepackage{enumitem}
\usepackage{titlesec}
\usepackage{xcolor}
\usepackage[hidelinks]{hyperref}
\usepackage{tabularx}

\hypersetup{
  pdftitle={Keshav Krishna - Cloud and ML Systems Resume},
  pdfauthor={Keshav Krishna},
  pdfsubject={Software and Machine Learning Engineer}
}

\pagestyle{empty}
\setlength{\parindent}{0pt}
\setlength{\tabcolsep}{0pt}
\setlength{\emergencystretch}{2em}
\urlstyle{same}

\titleformat{\section}
  {\large\scshape}
  {}{0pt}{}
  [\vspace{0.1em}\titlerule]
\titlespacing*{\section}{0pt}{1em}{0.55em}

\setlist[itemize]{
  label=--,
  leftmargin=1.8em,
  itemsep=0.25em,
  topsep=0.25em,
  parsep=0pt,
  partopsep=0pt
}

\newcommand{\entry}[4]{%
  \begin{tabularx}{\textwidth}{@{}Xr@{}}
    \textbf{#1} & #2 \\
    \textit{#3} & \textit{#4}
  \end{tabularx}\vspace{-0.2em}
}

\newcommand{\project}[2]{%
  \begin{tabularx}{\textwidth}{@{}Xr@{}}
    \textbf{#1} & \textit{#2}
  \end{tabularx}\vspace{-0.25em}
}

\begin{document}
\fontsize{10.1}{12.7}\selectfont

\begin{center}
  {\fontsize{24}{27}\selectfont\scshape Keshav Krishna}\\[0.12em]
  New York City, US $\mid$ +1 (479) 685 5463 $\mid$ \href{mailto:kk6081@nyu.edu}{kk6081@nyu.edu}\\[-0.05em]
  \href{https://github.com/k-keshav-k}{GitHub} $\mid$
  \href{https://k-keshav-k.github.io/}{Website} $\mid$
  \href{https://www.linkedin.com/in/keshav-krishna-5054681a9/}{LinkedIn} $\mid$
  \href{https://scholar.google.com/citations?user=zN7P_v4AAAAJ}{Google Scholar}\\[0.08em]
  \textbf{Graduation: May 2027 $\mid$ Available to start: May 2027}
\end{center}

\vspace{0.25em}
Software and machine learning engineer with three years of experience building production recommendation, personalization, and LLM systems for a consumer product serving 100M+ users. Experienced across the full ML lifecycle: data pipelines, model development, debugging, training, evaluation, inference, deployment, profiling, and iteration. Strong foundation in Python, C/C++, Linux, algorithms, system design, distributed data processing, and ML platform tooling.

\section{Education}

\entry{New York University, Courant Institute --- MS, Computer Science}{Aug 2025 -- May 2027}{GPA 3.95/4.0; Machine Learning, Building LLM Reasoners, Honors Algorithms}{New York City}

\vspace{0.4em}
\entry{Indian Institute of Technology, Ropar --- BTech, Computer Science and Engineering}{Aug 2019 -- Jun 2023}{GPA 8.74/10; Merit Scholarship for top academic performance}{Ropar, India}

\section{Experience}

\entry{Moore Capital Management --- Risk Technology Intern}{May 2026 -- Present}{Risk analytics and scalable data tooling; full-time summer, part-time fall}{New York City}
\begin{itemize}
  \item Built Python decision-support tools for portfolio and risk teams in a hedge-fund environment, combining data from multiple sources into clear, actionable analytics used in daily workflows.
  \item Developed and benchmarked distributed data-processing tools across large-scale compute environments; used profiling and system modeling to identify scalability bottlenecks and guide infrastructure decisions.
  \item Optimized Value at Risk (VaR), Conditional Value at Risk (CVaR), beta, and correlation calculations with vectorized Python and Numba JIT compilation; benchmarked implementations to reduce runtime on large financial datasets.
\end{itemize}

\vspace{0.5em}
\entry{JioSaavn --- Software Engineer, Data Science}{Jul 2023 -- Jul 2025}{Production recommendation, personalization, and language-model systems}{Mumbai, India}
\begin{itemize}
  \item Built recommendation and personalization systems with embeddings, vector retrieval, Spark, and Hive tables backed by Azure Blob Storage through ABFS, driving \textbf{1M impressions, 800K clicks, and 500K streams per day} in a 100M+ user product.
  \item Owned the ML lifecycle for a LLaMA3-8B-Instruct conversational playlist system, from intent modeling and retrieval-driven song selection through evaluation, inference, and deployment with Docker and Kubernetes.
  \item Designed a reproducible evaluation across SetFit, HingRoBERTa, mURIL, and LLaMA3-8B-Instruct for multilingual moderation; built a labeled test set and selected a model achieving \textbf{85\% accuracy}.
  \item Collaborated with product, editorial, and infrastructure teams to translate ambiguous requirements into production models, data pipelines, and services.
\end{itemize}

\vspace{0.5em}
\entry{GE Healthcare --- Software Intern, Medical Computer Vision}{May 2022 -- Jul 2022}{Detection and de-identification pipelines for clinical imaging}{Bangalore, India}
\begin{itemize}
  \item Built and evaluated YOLO-based medical-image de-identification and RGB/IR people-detection pipelines over 20K+ images.
\end{itemize}

\section{Selected ML Systems Projects}

\project{FlashAttention-2 --- Triton GPU Kernels \href{https://github.com/k-keshav-k/Flash-Attention}{[GitHub]}}{2026}
\begin{itemize}
  \item Implemented forward and backward attention kernels with tiled parallel work partitioning, online softmax, causal masking, and activation recomputation; reduced attention memory from $O(N^2)$ to $O(N)$ and enabled 65K-token sequences.
  \item Built a Linux GPU profiling and benchmarking suite for Transformer models up to 2.7B parameters using NVIDIA Nsight Systems, CUDA memory snapshots, NVTX, mixed precision, and \texttt{torch.compile} baselines.
\end{itemize}

\newpage
\section{Selected ML Systems Projects (continued)}

\project{Transformer Language Model and Training Stack from Scratch \href{https://github.com/k-keshav-k/Large-Language-Models-from-Scratch}{[GitHub]}}{2026}
\begin{itemize}
  \item Built a 17M-parameter PyTorch language model from low-level components: BPE tokenization, RoPE, RMSNorm, SwiGLU, causal attention, AdamW, gradient clipping, checkpointing, and cosine learning-rate scheduling.
  \item Implemented multiprocessing-based parallel tokenization, memory-mapped data loading, automated training sweeps, FLOPs and memory accounting, and systematic architecture ablations for reproducible model development and evaluation.
\end{itemize}

\vspace{0.55em}
\project{Pico-LLM --- Training, Evaluation, and Inference Platform \href{https://github.com/k-keshav-k/pico-llm}{[GitHub]}}{2025}
\begin{itemize}
  \item Developed a unified platform for SFT and DPO/GRPO training, checkpoint management, batch evaluation, and inference across Transformer, LSTM, and K-gram architectures.
  \item Added KV-cached generation, multi-GPU batch generation, automatic architecture detection, and robust token handling to improve inference throughput and experiment reliability.
\end{itemize}

\vspace{0.55em}
\project{Phase-Adaptive Generation for Diffusion LMs \href{https://github.com/k-keshav-k/Phase-Assisted-Generation}{[GitHub]}}{Feb 2026 -- May 2026}
\begin{itemize}
  \item Analyzed per-token refinement traces and trained Transformer predictors over 138K control tuples to assign dynamic per-block inference budgets.
  \item Integrated the predictor into the LLaDA decoding loop and matched AdaBlock's 89.5\% GSM8K accuracy with approximately \textbf{21\% fewer forward evaluations}; designed experiments to validate quality-efficiency tradeoffs. \href{https://github.com/k-keshav-k/Phase-Assisted-Generation/blob/main/writeup/final_report.pdf}{[Report]}
\end{itemize}

\vspace{0.55em}
\project{ML-Driven Hardware Resource Optimization \href{https://github.com/k-keshav-k/ML-based-hardware-optimization-using-Phases}{[GitHub]}}{Jan 2026 -- May 2026}
\begin{itemize}
  \item Modeled temporal traces from multicore workloads, clustered hardware pressure states, and distilled them into compact predictors with sub-KB model state for dynamic resource-management decisions.
\end{itemize}

\section{Research and Publications}

\begin{itemize}
  \item \href{https://www.frontiersin.org/journals/artificial-intelligence/articles/10.3389/frai.2025.1441250/full}{\textit{Advancements in Cache Management: A Review of Machine Learning Innovations}} --- Frontiers in AI, 2025. Single-author review of 40+ papers on ML-driven cache and resource management.
  \item \href{https://arxiv.org/abs/2410.15344}{\textit{LLC Intra-set Write Balancing for Non-Volatile Memory Caches}} --- arXiv, 2024. Bachelor's thesis research on system-level cache optimization.
\end{itemize}

\section{Technical Skills}

\textbf{Programming and systems:} Python, C/C++, JavaScript, Linux, data structures and algorithms, system design, \mbox{multiprocessing}, distributed systems, Git\\
\textbf{ML development lifecycle:} PyTorch, Hugging Face, vLLM, Triton, CUDA, model development, debugging, training, evaluation, inference, recommendation systems, vector search, SFT, GRPO, RLHF\\
\textbf{Platforms and infrastructure:} Azure (Blob Storage, ABFS), Google Cloud Platform (GCP), Kubernetes, Docker, Spark, Hive, Airflow, Redis, Qdrant, Slurm/HPC, NVIDIA Nsight Systems, Numba JIT, FastAPI, Postgres

\end{document}
